\documentclass[11pt,letterpaper]{article}
\usepackage[margin=1in]{geometry}\usepackage{amsmath,amssymb,booktabs,array,microtype}\usepackage[T1]{fontenc}\usepackage{lmodern}
\usepackage[colorlinks=true,linkcolor=blue,citecolor=blue,urlcolor=blue]{hyperref}
\setlength{\parindent}{0pt}\setlength{\parskip}{5pt}\sloppy
\title{A Clock That Must Run Fast:\\ the Solar System as an Inequality\\ on the Rate of Cosmic Time}
\author{Carl P.\ Zimmerman\\ \small Briar Creek Tech\\ \small AI-assisted research programme; not peer reviewed}
\date{12 September 2026}
\begin{document}\maketitle
\begin{abstract}
A cuscuton-clock sector with a projected-gradient scalar, coupled to matter disformally along the clock's own unit timelike direction, gives $\gamma_{\rm PPN}=1$ exactly and no added vector field \cite{paper24}. Boosted, the same coupling gives $\alpha_1=8f_s$, five orders over the bound unless local matter drags the clock into its own rest frame to within $4.6\,$m$\,$s$^{-1}$. That alignment was left uncomputed. We compute it. Four structural results precede the number. \emph{(i)} Writing the matter metric with independent conformal and disformal strengths $A=1-2a\varphi$, $B=-2b\varphi$ and demanding $\gamma_{\rm PPN}=1$ forces $b=2a$ exactly: the operator that fixes light bending \emph{is} the operator that generates the preferred-frame effect, so the gate cannot be evaded by tuning. \emph{(ii)} The cuscuton's linearised Lagrangian has a vanishing coefficient for $\dot\psi^2$, so the clock propagates nothing and its tilt solves an elliptic boundary-value problem whose condition at infinity is the cosmic frame. \emph{(iii)} The stiffness resisting the drag is the gradient function at the \emph{local} invariant, $W=f_s^2GM^2/(8\pi s_0r^4)$, which is free of both the matter coupling $\lambda$ and the MOND coefficient $\beta$: neither can buy alignment, and the clock rate $s_0$ is the only handle. \emph{(iv)} The gain at a stellar surface is $D(R)=24s_0/f_s$, independent of the star's mass and radius entirely, and the falling stiffness makes the exterior tilt decay as $r^{-(\sqrt{17}-1)/2}=r^{-1.562}$ rather than the $r^{-3}$ of a constant-stiffness dipole. Both help; neither suffices. Carried to Earth's orbit the drag is $0.011$ at the minimum clock rate and the residual misses the bound by $7.9\times10^4$. Inverting, the gate becomes $s_0\gtrsim1.5\times10^7$. \textbf{The solar system does not exclude this construction; it demands a clock running about ten million times proper time.} That is the third independent constraint on $s_0$ and the third pointing the same way, after the clock identity and positivity of the sector's energy --- but the first two want order unity, and nothing in the programme explains a number seven orders larger. It is recorded as a requirement, not a derivation. The price, also computed: the margin falls to $m_{\rm rel}=6.8\times10^{-12}$ and the scalar's kinetic coefficient rises to $P_X/d=1.5\times10^{11}$, a strong-coupling question this paper states and does not answer. Four new Lean~4 certificates; the file stands at 128 theorems, zero \texttt{sorry}.
\end{abstract}

\section{The question handed on}
In \cite{paper24} the dark sector of \cite{paper13} was placed on the locus $W_0\equiv U-2\gamma\bar q^2\dot{\bar q}=0$, where the clock--scalar mixing vanishes identically, and matter was coupled disformally along the clock's own unit timelike direction $n_\mu=-\partial_\mu\tau/s$. Expanding $\tilde g=(1-2\varphi)g-4\varphi\,n\otimes n$ in the weak field gives $\tilde\Phi=\Phi+\varphi$ and $\tilde\Psi=\Psi+\varphi$: the shift is common to both potentials, so it cancels from their difference and $\gamma_{\rm PPN}=1$ exactly, however strong the scalar. That is the TeVeS mechanism \cite{teves} obtained without adding a vector field, because the cuscuton \cite{cuscuton} already carries the direction and propagates no scalar of its own.

For a system moving at $w$ through the clock's frame the same term gives $\tilde g_{0i}=4\varphi w_i$, first order in the velocity and first order in the scalar, which is the preferred-frame signature. Matching to the post-Newtonian form gives $\alpha_1=8f_s$ with $f_s=\varphi/\Phi_N$ the scalar's share of the local potential. At the share MOND requires at solar-system accelerations, $f_s\approx1$, that is $8$ against a bound of $10^{-4}$. The bound converts into an alignment requirement of $4.6\,$m$\,$s$^{-1}$ out of the solar system's $370\,$km$\,$s$^{-1}$, one part in $8\times10^4$. Whether the clock's own field equation delivers that is the subject of this paper.

\section{The coupling cannot be switched off}
Write the matter metric with independent strengths, $\tilde g_{\mu\nu}=A(\varphi)g_{\mu\nu}+B(\varphi)n_\mu n_\nu$ with $A=1-2a\varphi$ and $B=-2b\varphi$. Using $n_0^2=1+2\Phi$,
\begin{equation}
\tilde\Phi=\Phi+(b-a)\varphi,\qquad \tilde\Psi=\Psi+a\varphi ,
\end{equation}
so that with no anisotropic stress in the Einstein frame, $\Psi=\Phi$,
\begin{equation}
\gamma_{\rm PPN}=\frac{\Phi+a\varphi}{\Phi+(b-a)\varphi}=1\qquad\Longleftrightarrow\qquad \boxed{b=2a}.
\end{equation}
The disformal strength is locked to the conformal one. There is no limit in which light bending stays general-relativistic and the preferred-frame source is small: they are the same operator at a fixed relative weight. The gate must be met by dragging or not at all. (Lean: \texttt{disformal\_coefficient\_locked}.)

\section{The clock is a constraint, so the tilt is a boundary-value problem}
Write $\tau=t+\psi(x,t)$. Then $s=\sqrt{(1+\dot\psi)^2-|\nabla\psi|^2}$, and expanding,
\begin{equation}
\text{coefficient of }\dot\psi^2\ \text{in } s\;=\;0,\qquad \text{coefficient of }|\nabla\psi|^2\;=\;-\tfrac12 .
\end{equation}
The linear-in-$s$ structure kills the time-kinetic term exactly and leaves $-\tfrac12W|\nabla\psi|^2$. The clock therefore carries no propagating degree of freedom and obeys a constraint on each slice. Varying the disformal point-mass action for a source of density $\rho$ moving at $v$,
\begin{equation}
\nabla\!\cdot\!\big(W\nabla\psi\big)=-4\,\nabla\!\cdot\!\big(\varphi\rho\,v\big).
\label{elliptic}
\end{equation}
This is elliptic, with the cosmic frame as its boundary condition at infinity. Whether the Sun wins is a straight competition between a local source and a stiffness --- there is no propagating response that may or may not arrive.

\section{The stiffness cannot be bought}
The coefficient in (\ref{elliptic}) is $W$ evaluated at the \emph{local} gradient invariant, not at $W(0)=U$, because $sW(Y,\tau)$ expands as $W-\tfrac12W|\nabla\psi|^2$ with $W$ taken where the field point is. In the saturated high-acceleration regime $W\simeq kY$ with $k=\lambda^2/(8\pi Gs_0)$ \cite{paper24}, and the scalar gradient that gives matter a share $f_s$ of the Newtonian acceleration is $|\nabla\varphi|=f_sGM/(\lambda r^2)$. Hence
\begin{equation}
W_{\rm local}=\frac{f_s^2GM^2}{8\pi s_0r^4},\qquad \frac{\partial W_{\rm local}}{\partial\lambda}=\frac{\partial W_{\rm local}}{\partial\beta}=0 .
\end{equation}
\emph{Both free coefficients cancel.} Neither the matter coupling nor the MOND coefficient can be chosen to improve the alignment. The clock rate $s_0$ is the only remaining handle, which is what makes the gate an inequality on a single quantity.

\section{The gain, and how far it reaches}
At a stellar surface, balancing source against stiffness in (\ref{elliptic}) with $\varphi_s=f_sGM/R$ and $\rho_\odot=3M/(4\pi R^3)$,
\begin{equation}
D(R)\;=\;\frac{4\varphi_s\rho_\odot}{W_{\rm local}(R)}\;=\;\frac{24\,s_0}{f_s},\qquad \frac{\partial D}{\partial M}=\frac{\partial D}{\partial R}=0 .
\end{equation}
The gain is independent of the star entirely --- a property of the theory, not of the object. At the clock rate that positivity already forces, $s_0\ge2$ \cite{paper24}, it is at least $48$: the source beats the stiffness and the clock \emph{is} dragged at the surface. (Lean: \texttt{drag\_gain\_is\_star\_independent}.)

Outside, with $W\propto r^{-4}$, the $\ell=1$ mode of (\ref{elliptic}) obeys $r^2g''-2rg'-2g=0$, indicial equation $n^2-3n-2=0$, roots $(3\pm\sqrt{17})/2$. The decaying branch gives a tilt falling as
\begin{equation}
|\nabla\psi|\;\propto\;r^{-(\sqrt{17}-1)/2}\;=\;r^{-1.562},
\end{equation}
far slower than the $r^{-3}$ of a constant-stiffness dipole. The clock is dragged much further out than a naive estimate would give, and this works in the construction's favour. (Lean: \texttt{drag\_exponent\_root}.)

\section{The number}
Carrying the surface gain to Earth's orbit with that exponent, at $f_s=1$ and $s_0=2$,
\begin{equation}
D(1\,\text{AU})=0.011,\qquad \text{residual}=\frac{1}{1+D}=0.989 ,
\end{equation}
against a requirement of $1.25\times10^{-5}$: short by $7.9\times10^4$. Since $D$ is linear in $s_0$, the bound inverts (Lean: \texttt{alignment\_forces\_clock\_rate}) into
\begin{equation}
\boxed{\;s_0\;\gtrsim\;1.5\times10^{7}\;}
\end{equation}
\textbf{The solar system does not exclude this construction. It demands a clock running about ten million times proper time.}

\begin{center}
\begin{tabular}{lll}
\toprule
constraint on $s_0$ & source & value \\
\midrule
clock identity $s_0-1=w/m_{\rm rel}$ & conservation \cite{paper21} & $s_0>1$ \\
positivity of the sector's energy & \cite{paper24} & $s_0\ge2$ \\
solar-system alignment & this paper & $s_0\gtrsim1.5\times10^7$ \\
\bottomrule
\end{tabular}
\end{center}

Three independent routes, all pushing the same quantity the same way. That convergence is real and is the useful part. But the first two want order unity and the third wants ten million, and nothing in this programme explains a number of that size. It is recorded as a requirement, not a derivation.

\section{What it costs elsewhere}
At the demanded rate, $m_{\rm rel}=w/(s_0-1)=6.8\times10^{-12}$, the cubic operator's share of the sector density falls to $6.8\times10^{-8}$, and the scalar's kinetic coefficient rises to
\begin{equation}
\frac{P_X}{d}=\frac{1}{m_{\rm rel}}=1.5\times10^{11}.
\end{equation}
Nothing computed contradicts this. The sound speed $c_s^2=-w/(2-m_{\rm rel})$ and every gate of \cite{paper24} depend on $w$ and not on $s_0$, so they are unchanged. But a kinetic coefficient eleven orders above its neighbour is a strong-coupling question, and this paper states it rather than answering it.

\section{What is not claimed}
\textbf{This is not a complete theory of gravity and nothing here should be read as one.} The interior match of \S5 evaluates source and stiffness at the stellar surface with a uniform-density star; the real interior profile of the gradient invariant is not solved and the gain could move by an order in either direction. The exterior stiffness is taken as $r^{-4}$ throughout, which assumes the scalar's share $f_s$ stays constant from the surface to $1\,$AU. The elliptic problem is solved for its indicial exponents and matched at one radius, not integrated through the star. The post-Newtonian matching inherited from \cite{paper24} is convention-dependent at the factor-of-two level, so the coefficient $8$ is uncertain by that much and the five-orders conclusion is not. The bound used is $|\alpha_1|<10^{-4}$; a tighter bound raises $s_0$ proportionally. The origin of $s_0\sim10^7$ is not explained, the strong-coupling question of \S7 is not answered, and $\kappa=\tfrac12$ remains fitted and provably underivable by this class of actions \cite{repo}.

\section{Reproducibility}
\texttt{fable\_independent\_2026/L216\_clock\_alignment.py} with its committed output and results file (commit \texttt{bc15f3e24}). The Lean file \texttt{fable\_independent\_2026/lean\_2026/Mondlean.lean} carries \texttt{disformal\_coefficient\_locked}, \texttt{drag\_gain\_is\_star\_independent}, \texttt{drag\_exponent\_root} and \texttt{alignment\_forces\_clock\_rate}; 128 theorems, exit 0, zero \texttt{sorry}, axioms $\subseteq\{$\texttt{propext}, \texttt{Classical.choice}, \texttt{Quot.sound}$\}$.

\begin{thebibliography}{9}
\bibitem{repo} The research repository, \url{https://github.com/carlzimmerman/zimmerman-formula} (2026).
\bibitem{paper13} C.~P.~Zimmerman, \emph{A Primordial-Clock Relativistic MOND Candidate: Status Report}, doi:10.5281/zenodo.22699828 (2026).
\bibitem{paper21} C.~P.~Zimmerman, \emph{The Clock Rate Is the Dark Matter Pressure}, doi:10.5281/zenodo.22728545 (2026).
\bibitem{paper24} C.~P.~Zimmerman, \emph{The Decoupling Locus: a MOND Scale Derived from a Clock, and the Solar System Reduced to One Number}, doi:10.5281/zenodo.22729429 (2026).
\bibitem{teves} J.~D.~Bekenstein, Phys.\ Rev.\ D \textbf{70}, 083509 (2004).
\bibitem{aest} C.~Skordis, T.~Z\l{}o\'snik, Phys.\ Rev.\ Lett.\ \textbf{127}, 161302 (2021).
\bibitem{cuscuton} N.~Afshordi, D.~J.~H.~Chung, G.~Geshnizjani, Phys.\ Rev.\ D \textbf{75}, 083513 (2007).
\bibitem{will} C.~M.~Will, \emph{Theory and Experiment in Gravitational Physics}, 2nd ed., Cambridge (2018).
\end{thebibliography}
\end{document}
